\scanpage{83}
Les $\alpha_{p,q}$ permettent de définir des
\uline{\unsure{Définissons|Définitions}}.
\[
P_{\unsure{n|p}}(X,Y)=\operatorname{Hom}_{P_1}(\Gamma^n X,Y)
\]
\inserted{donc
\[
\left\{\begin{aligned}
P_1(X,Y)&=\operatorname{Hom}_{P_1}(X,Y)\\
P_0(X,Y)&=\operatorname{Hom}_{P_1}(\mathbb{1},Y)
\end{aligned}\right.
\]}
et des applications de composition
\[
\begin{aligned}
P_p(X,Y)\times P_q(Y,Z)&\longrightarrow P_{pq}(X,Z)\\
f\qquad g&\longmapsto g\circ f
\end{aligned}
\]
par
\diagram{\begin{tikzcd}[ampersand replacement=\&,column sep=small]
\Gamma^p X\arrow[r,"f"] \& Y \& \Gamma^qY\arrow[r,"g"] \& Z\\
{\Gamma^q(\Gamma^p(X))}\arrow[urr,"{\Gamma^q(f)}"]\arrow[r,dash,"\deleted{\unsure{\Gamma^p}}"'] \& \deleted{\unsure{\Gamma^q(Y)}}\arrow[rr,dash] \& \& T\\
{\Gamma^{pq}(X)}\arrow[u,"{\alpha_{q,p}^X}"] \& \& \& {g\circ f=g\Gamma^q(f)\alpha_{q,p}^X}
\end{tikzcd}}

Le diagramme comm. $(*)$ assure l'associativité
de la composition, $(**)$ \unsure{tout uniment} exprime que si \deleted{$g$ \unsure{resp.} $f$ est
de degré $1$,} \inserted{$p$ ou $q=1$,} alors le composé est celui
provenant de la bifonctorialité de
$P_p(X,Y)$ sur $P_1^\circ\times P_1$, enfin les compatibilités
$(***)$ signifient, pour la \unsure{deuxième|dernière|dernier}, que
\[
g\circ f=g\qquad\text{si $g$ de degré $0$}
\]
la deuxième signifie
\diagram{\begin{tikzcd}[ampersand replacement=\&]
|[draw,ellipse,inner sep=6pt]|{\overset{\deleted{\unsure{\Gamma}}}{\mathbb{1}}\xrightarrow{\substack{\alpha_{q,0}\\\sim}}\Gamma^q(\mathbb{1})}
\end{tikzcd}}
% slanted
\margin{\unsure{plus tard|pour plus tard}}
\unsure{on|je} l'avait oublié dans les axiomes
préliminaires -- il ne semble pas
\inserted{pas se déduire du \unsure{reste|restant}\unsure{.|\ldots}}
% archive sheet number: 82
